% GitHub/reducesymm/latticeQED/CviVat26.tex   compile by  pdflatex CviVat26; biber CviVat26

% arara: pdflatex : { synctex: yes }
% arara: biber
% arara: pdflatex : { synctex: yes }

%%%%%%%%%%%%%%%%%%%%%%% logical setup, no need to edit %%%%%%%%%%%%%
\newif\ifboyscout             %%
\boyscouttrue 					%% commented, internal version %%

%%%%% Toggle between draft and non-draft versions %%%%%%%%%%%%%%%%%%%%%%%
%\boyscoutfalse                 % public, for hyperlinked CviVat26.pdf
%%%%%%%%%%%%%%%%%%%%%%%%%%%%%%%%%%%%%%%%%%%%%%%%%%%

\documentclass[11pt,a4paper]{article}
\usepackage[utf8]{inputenc}
\usepackage[T1]{fontenc}
\usepackage{amsmath,amssymb,amsfonts}
\usepackage{booktabs}
\usepackage{geometry}
\geometry{margin=2.5cm}
\usepackage{graphicx}
\usepackage{caption}
% defsCviVat26.tex

%  Predrag          				11aug2026
% copied from dasbuch/book/inputs/def.tex  2025-02-08  version

%%%%%%%%%%%%%%%%%%%%%%%%%%%%%%%%%%%%%%%%%%%%%%%%%%%%%%%%%%%%%%%%%%%%%%%%%
%% defines macros used throughout ChaosBook and related
%%%%%%%%%%%%%%%%%%%%%%%%%%%%%%%%%%%%%%%%%%%%%%%%%%%%%%%%%%%%%%%%%%%%%%%%%

%%%%%%%%%%%%%%% EQUATIONS %%%%%%%%%%%%%%%%%%%%%%%%%%%%%%%
\newcommand{\beq}{\begin{equation}}
	\newcommand{\continue}{\nonumber \\ }
	\newcommand{\nnu}{\nonumber}
	\newcommand{\eeq}{\end{equation}}
\newcommand{\ee}[1] {\label{#1} \end{equation}}
\newcommand{\bea}{\begin{eqnarray}}
\newcommand{\ceq}{\nonumber \\ & & }
\newcommand{\eea}{\end{eqnarray}}
\newcommand{\barr}{\begin{array}}
\newcommand{\earr}{\end{array}}

%%%%%%%%%%%%%%% REFERENCING EQUATIONS ETC. %%%%%%%%%%%%%%%%%%%%%%%%%%%%%%%
\newcommand{\rf}     [1] {~\cite{#1}}
\newcommand{\refref} [1] {ref.~\cite{#1}}
\newcommand{\refRef} [1] {Ref.~\cite{#1}}
\newcommand{\refrefs}[1] {refs.~\cite{#1}}
\newcommand{\refRefs}[1] {Refs.~\cite{#1}}
\newcommand{\refeq}  [1] {(\ref{#1})}
% in amstex, \eqref is predefined and better than \refeq
\newcommand{\refeqs} [2]{(\ref{#1}--\ref{#2})}
\newcommand{\refpage}[1] {page~\pageref{#1}}
\newcommand{\reffig} [1] {figure~\ref{#1}}
\newcommand{\reffigs} [2] {figures~\ref{#1} and~\ref{#2}}
\newcommand{\refFig} [1] {Figure~\ref{#1}}
\newcommand{\refFigs} [2] {Figures~\ref{#1} and~\ref{#2}}
\newcommand{\reftab} [1] {table~\ref{#1}}
\newcommand{\refTab} [1] {Table~\ref{#1}}
\newcommand{\reftabs}[2] {tables~\ref{#1} and~\ref{#2}}
\newcommand{\refsect}[1] {sect.~\ref{#1}}
\newcommand{\refsects}[2] {sects.~\ref{#1} and \ref{#2}}
\newcommand{\refSect}[1] {Sect.~\ref{#1}}
\newcommand{\refSects}[2] {Sects.~\ref{#1} and \ref{#2}}
\newcommand{\refsecttosect}[2] {Sects.~\ref{#1} to~\ref{#2}}
%\newcommand{\refchap}[1] {chapter~\ref{#1}}
%\newcommand{\refChap}[1] {Chapter~\ref{#1}}
%\newcommand{\refchaps}[2] {chapters~\ref{#1} and~\ref{#2}}
%\newcommand{\refchaptochap}[2] {chapters~\ref{#1} to~\ref{#2}}
\newcommand{\refappe}[1] {appendix~\ref{#1}}
\newcommand{\refappes}[2] {appendices~\ref{#1} and~\ref{#2}}
\newcommand{\refAppe}[1] {Appendix~\ref{#1}}
\newcommand{\refrem} [1] {remark~\ref{#1}}
\newcommand{\refRem} [1] {Remark~\ref{#1}}
%\newcommand{\refexam}[1] {example~\ref{#1}}
%\newcommand{\refExam}[1] {Example~\ref{#1}}
%\newcommand{\refexer}[1] {exercise~\ref{#1}}
%\newcommand{\refExer}[1] {Exercise~\ref{#1}}
%\newcommand{\refsolu}[1] {solution~\ref{#1}}
%\newcommand{\refSolu}[1] {Solution~\ref{#1}}

%%%%%%%% DISPLAY COMMENTS IN THE TEXT %%%%%%%%%%%%%%%%%%%%
\ifboyscout 
% also search the text for lines starting with %%  to
% locate various internal comments, recent edits etc.
\typeout{============ COMMENTED =====}
\newcommand{\PC}[2]{\begin{quote}{\color{red}[#1 Predrag] \small #2}\end{quote}}
\newcommand{\PCedit}[1]{{\color{green}#1}}
\newcommand{\GV}[2]%{$\footnotemark\footnotetext{G{\'a}bor: #1}$}
{\begin{quote}\GVedit{[#1 G{\'a}bor] \small #2}\end{quote}}
\newcommand{\GVedit}[1]{{\color{blue}#1}}
\newcommand{\Private}[1]{{\color{blue}#1}}
\newcommand{\toCB}{\marginpar{\footnotesize 2CB}}  % to compare with ChaosBook
\newcommand{\inCB}{\marginpar{\footnotesize now in CB}} % entered in ChaosBook
\newcommand{\CBlibrary}[1]
{\href{https://ChaosBook.org/library/#1.pdf} { (click here)}}
\else % drop comments
% do not turn on labeling of equations on margins
\typeout{============ UNCOMMENTED =====}
\newcommand{\PC}[2]{}{}
\newcommand{\PCedit}[1]{#1}
\newcommand{\GV}[2]{}
\newcommand{\GVedit}[1]{#1}
\newcommand{\Private}[1]{}
\newcommand{\toCB}{}
\newcommand{\inCB}{}
\newcommand{\CBlibrary}[1]{}
\fi  %%%%%%%%%%%% END OF ON/OFF COMMENTS SWITCH %%%%%%%%%%%%%%%%%%%%

%%%%%%%%%%%%%%%  prepare hyperlinked pdf %%%%%%%%%%%%%%% 
\newcommand{\wwwcb}[1]{       % keep homepage flexible:
	{\href{http://ChaosBook.org#1}
		{ChaosBook.org#1}}}
\newcommand{\wwwgt}{{\href{http://birdtracks.eu}
		{birdtracks.eu}}}
\newcommand{\wwwQFT}[1]{
	{\href{http://ChaosBook.org/FieldTheory#1}
		{ChaosBook.org/\-Field\-Theory#1}}}
\newcommand{\wwwcnsQFT}[1]{
	{\href{http://ChaosBook.org/FieldTheory#1}
		{ChaosBook.org/\-Field\-Theory#1}}}
\newcommand{\weblink}[1]{{\href{http://#1}{#1}}}
\newcommand{\HREF}[2]{
	{\href{#1}{#2}}}
\newcommand{\mpArc}[1]{
	{\href{http://www.ma.utexas.edu/mp_arc-bin/mpa?yn=#1}
		{mp\_arc~#1}}}
\newcommand{\arXiv}[1]{
	{\href{https://arXiv.org/abs/#1}{arXiv:#1}}}
\newcommand{\bioRxiv}[1]{
	{\href{https://biorxiv.org/content/10.1101/#1}{bioRxiv:#1}}}
	
% online PDF \toChaosBook{}{} section pointer       % Predrag 2019-12-07
%                        %#1 = PDF file section, #2 = section title
\newcommand{\toChaosBook}[2]
{\HREF{http://ChaosBook.org/chapters/ChaosBook.pdf\##1}{ChaosBook\,#2}}
\newcommand{\toCLpdf}[2]
{\HREF{https://arxiv.org/pdf/2503.22972\##1}{CL18\, #2}}
% if edited, then: {\HREF{https://ChaosBook.org/overheads/spatiotemporal/CL18.pdf\##1}{CL18\, #2}}

% online PDF \toBirdtracks{}{} section pointer       % Predrag 2020-01-02
%                        %#1 = PDF file section, #2 = section title
\newcommand{\toBirdtracks}[2]
{\HREF{http://birdtracks.eu/version9.0/GroupTheory.pdf\##1}{birdtracks\,#2}}

    %%%%%%% Number, group theory %%%%%%%%%%%%%%%%%%%%%%%%%%%%
\newcommand{\reals}{\mathbb{R}}
\newcommand{\complex}{\mathbb{C}}
\newcommand{\integers}{\mathbb{Z}}
\newcommand{\rationals}{\mathbb{Q}}
\newcommand{\naturals}{\mathbb{N}}
\renewcommand\Im{\ensuremath{\mathrm{Im}\,}}
\renewcommand\Re{\ensuremath{\mathrm{Re}\,}}
\renewcommand{\det}{\mathrm{det}\,}
\newcommand{\Det}{\mathrm{Det}\,}
\newcommand{\tr}{\mathrm{tr}\,}
\newcommand{\Tr}{\mathrm{Tr}\,}
\newcommand{\sign}[1]{\sigma_{#1}}
%\newcommand{\sign}[1]{{\rm sign}(#1)}
\newcommand{\shift}{\ensuremath{\mathsf{r}}}        % in LC21
\newcommand{\Un}[1]{\ensuremath{\textrm{U}(#1)}}         % in DasBuch
\newcommand{\SUn}[1]{\ensuremath{\textrm{SU}(#1)}}         % in DasBuch
%\newcommand{\On}[1]{\ensuremath{\mathbf{O}(#1)}}
\newcommand{\On}[1]{\ensuremath{\textrm{O}(#1)}}
%\newcommand{\SOn}[1]{\ensuremath{\mathbf{SO}(#1)}} % in Siminos thesis
\newcommand{\SOn}[1]{\ensuremath{\textrm{SO}(#1)}}         % in DasBuch
\newcommand{\Spn}[1]{\ensuremath{\textrm{Sp}(#1)}}         % in DasBuch
\newcommand{\Cn}[1]{\ensuremath{\textrm{C}_{#1}}}
\newcommand{\Cnv}[1]{\ensuremath{\textrm{C}_{#1v}}}
%\newcommand{\Dn}[1]{\ensuremath{\mathbf{D}_{#1}}    % in Siminos thesis
\newcommand{\Dn}[1]{\ensuremath{\textrm{D}_{#1}}}              % in DasBuch
\newcommand{\Zn}[1]{\ensuremath{\textrm{Z}_{#1}}}   % C_n in DasBuch
\newcommand{\SLn}[2]{\ensuremath{\textrm{SL}(#1,#2)}} % wiki convention
% \newcommand{\SLn}[2]{\ensuremath{\textrm{SL}_{#1}(#2)}} % another convention

     %%%%%%%%%% periods: %%%%%%%%%%%%%%%%%%%%%%%%%%%%
\newcommand\period[1]{{\ensuremath{T_{#1}}}}         %continuous cycle period
%\newcommand\period[1]{{\tau_{#1}}}
\newcommand{\cl}[1]{{\ensuremath{n_{#1}}}}   % discrete length of a cycle, Predrag
%\newcommand{\cl}[1]{|#1|}  % the length of a periodic orbit, Ronnie
\newcommand{\speriod}[1]{{\ensuremath{L_{#1}}}}  %spatial period
\newcommand{\tilt}[1]{{\ensuremath{S_{#1}}}}  %relative shift

    %%%%%%% CL18 specific %%%%%%%%%%%%%%%%%%%%%%%%%%%%
\newcommand{\dft}{det\-er\-min\-is\-tic field theory}
\newcommand{\Dft}{Det\-er\-min\-is\-tic field theory}
\newcommand{\dfts}{det\-er\-min\-is\-tic field theories}
\newcommand{\Dfts}{Det\-er\-min\-is\-tic field theories}
\newcommand{\detZeta}{det\-er\-min\-is\-tic zeta func\-tion}
\newcommand{\DetZeta}{Det\-er\-min\-is\-tic zeta func\-tion}
\newcommand{\detSum}{det\-er\-min\-is\-tic part\-it\-ion sum}
\newcommand{\DetSum}{Det\-er\-min\-is\-tic part\-it\-ion sum}
\newcommand{\fundPip}{fundamental parallelepiped}
\newcommand{\brick}{block}
\newcommand{\Brick}{Block}
\newcommand{\jacobianOrb}{orbit Jacobian matrix}
\newcommand{\JacobianOrb}{Orbit Jacobian matrix}
\newcommand{\jacobianOrbs}{orbit Jacobian matrices}
\newcommand{\JacobianOrbs}{Orbit Jacobian matrices}
\newcommand{\jacobianOp}{orbit Jacobian operator}
\newcommand{\JacobianOp}{Orbit Jacobian operator}
\newcommand{\jacobianOps}{orbit Jacobian operators}
\newcommand{\JacobianOps}{Orbit Jacobian operators}
\newcommand{\Hilldet}{orbit Jacobian}
\newcommand{\HillDet}{Orbit Jacobian}
%\newcommand{\Hilldet}{Hill determinant}
%\newcommand{\HillDet}{Hill determinant}
\newcommand{\stabexp}{stability exponent}
\newcommand{\Stabexp}{Stability exponent}
\newcommand{\jMorb}{{\ensuremath{\cal \jMps}}}
\newcommand{\conf}{\ensuremath{x}}      % Configuration space coordinate
\newcommand{\lattice}{\ensuremath{\mathcal{L}}}
% \newcommand{\action}{S}
\newcommand{\ELe}{defining equations}           % CL18 definition
\newcommand{\ELE}{Defining equations}           % CL18 definition
% \newcommand{\ELe}{The Law}                    % elements definition
% \newcommand{\ELe}{Euler-Lagrange equation}    % initial CL18 definition
\newcommand{\jMat}{\mathbb{\jMps}}
\newcommand{\spt}{spa\-tio\-temp\-or\-al}
\newcommand{\Spt}{Spa\-tio\-temp\-or\-al}
\newcommand{\statesp}{state space}
\newcommand{\Statesp}{State space}
\newcommand{\stateDsp}{state-space}
\newcommand{\StateDsp}{State-space}
\newcommand{\po}{periodic orbit}
\newcommand{\Po}{Periodic orbit}
\newcommand{\rpo}{rela\-ti\-ve periodic orbit}
\newcommand{\Rpo}{Rela\-ti\-ve periodic orbit}
\newcommand{\ppo}{pre-periodic orbit}
\newcommand{\Ppo}{Pre-periodic orbit}
\newcommand{\eqv}{equi\-lib\-rium}
\newcommand{\Eqv}{Equi\-lib\-rium}
\newcommand{\eqva}{equi\-lib\-ria}
\newcommand{\Eqva}{Equi\-lib\-ria}
\newcommand{\reqv}{rela\-ti\-ve equi\-lib\-rium}
%   \newcommand{\reqv}{travelling wave}
\newcommand{\Reqv}{Rela\-ti\-ve equi\-lib\-rium}
%   \newcommand{\Reqv}{travelling wave}
\newcommand{\reqva}{rela\-ti\-ve equi\-lib\-ria}
\newcommand{\Reqva}{Rela\-ti\-ve equi\-lib\-ria}
\newcommand{\equilibrium}{equi\-lib\-rium}
\newcommand{\equilibria}{equi\-lib\-ria}
\newcommand{\Equilibria}{Equi\-lib\-ria}
\newcommand{\steady}{steady state}
\newcommand{\Steady}{Steady state}
\newcommand{\lst}{periodic state}
\newcommand{\Lst}{Periodic state}
\newcommand{\lsts}{periodic states}
\newcommand{\Lsts}{Periodic states}
 \newcommand{\primeOrb}{prime orbit}
\newcommand{\PrimeOrb}{Prime orbit}
\newcommand{\primeOrbs}{prime orbits}
\newcommand{\PrimeOrbs}{Prime orbits}
\newcommand{\fconf}{field con\-fig\-ura\-tion}
\newcommand{\fconfs}{field con\-fig\-ura\-tions}
\newcommand{\Fconf}{Field con\-fig\-ura\-tion}
\newcommand{\Fconfs}{Field con\-fig\-ura\-tions}
\newcommand{\ssp}{\ensuremath{\phi}}             % kittens lattice site field
\newcommand{\cssp}{\ensuremath{\tilde{\phi}}}    % Complex state space point
\newcommand{\Mm}{\ensuremath{\mathsf{M}}}
\newcommand{\Xx}{\ensuremath{\mathsf{\Phi}}}      % kittens lattice field
\newcommand{\xx}{\ensuremath{\mathsf{\bphi}}}      % kittens lattice field
\newcommand{\Source}{\mathsf{J}}
\newcommand{\source}{{j}}
\newcommand{\Aa}{\ensuremath{\bar{\A}}}
\newcommand{\A}{\ensuremath{\mathcal{A}}}       % alphabet
\newcommand{\bcs}{boundary conditions}
\newcommand{\pcell}{primitive cell}      % 2023-01-13 Standard name
\newcommand{\Pcell}{Primitive cell}      % 2023-01-13 Standard name
\newcommand{\Bcell}[1]{\ensuremath{\mathbb{#1}}} % 2023-02-13 `Bravais' cell
\newcommand{\Blatt}{Bravais lattice}
\newcommand{\BravCell}[3]{{\ensuremath{[#1\!\times\!#2]_{#3}}}}
\newcommand{\LTS}[3]{{\ensuremath{[\speriod{#1}\!\times\!\period{#2}]_\tilt{#3}}}}
\newcommand{\dmn}{-dim\-en\-sion\-al}  %  experimental 220ct2009
%\newcommand{\dmn}{\ensuremath{d}}  %  n-dimensional
%\newcommand{\dmn}{\ensuremath{\!-\!d}}  %  n-dimensional
\newcommand{\vol}{{V}}                  % volume of i-th tile
\newcommand{\wavenum}{\ensuremath{\mathsf{k}}}        % in LC18


%%%%%%%%%%%%%% Specific PDEs %%%%%%%%%%%%%%%%%%%%%%%%%%%%
\newcommand{\catlatt}{spa\-tio\-temp\-or\-al cat}     % Predrag, experimental
\newcommand{\catLatt}{Spa\-tio\-temp\-or\-al cat}     % Predrag, experimental
\newcommand{\PV}{Percival-Vivaldi}
\newcommand{\AW}{Adler-Weiss}
% \newcommand{\GO}{Gutkin-Osipov}
\newcommand{\KS}{Ku\-ra\-mo\-to-Siva\-shin\-sky}
\newcommand{\KSe}{Ku\-ra\-mo\-to-Siva\-shin\-sky equation}
\newcommand{\NS}{Navier-Stokes}
\newcommand{\pCf}{plane Couette flow}
\newcommand{\PCf}{Plane Couette flow}





	
%%%%%%%%%%%% REMOVE THIS EVENTUALLY %%%%%%%%%%%
   					%% macros, DasBuch style
\input{biblatex}
\addbibresource{../bibtex/QFT.bib}
\usepackage[pdftex,colorlinks]{hyperref}

\title{\Dft\ of lattice QED}
% 1st version, Gemini:  {Deterministic spatiotemporal evaluation of compact $U(1)$ lattice gauge theory: {Prime orbits} and stability spectra on the minimal primitive cell}
\author{Predrag Cvitanovi\'{c}, Gemini the Robot and her Master G{\'a}bor Vattay}
% 1st version: {Unsung Hero}
\date{\today}

\begin{document}
	
\maketitle
\tableofcontents

\begin{abstract}
We present a self-contained, linear mathematical derivation of the foundational classical prime orbits and their linear stability spectra for compact $U(1)$ lattice electrodynamics on the minimal Euclidean primitive cell $[2^3 \times 2]_0$. By formulating the deterministic Euler-Lagrange equations of motion and eliminating local gauge redundancy algebraically through the Maximal Spanning Tree gauge, we reduce the dynamical system to a non-singular 49-dimensional active coordinate subspace. In the pure gauge limit, we construct explicit analytical solutions for the trivial vacuum, constant-flux torons, and saddle-point sphalerons, alongside numerically relaxed antipodal monopole-antimonopole dipole defects. We evaluate the Morse index for each stationary state from the spectrum of the reduced {\jacobianOrb}, establishing the exact statistical weight contributions required for a deterministic spatiotemporal zeta function evaluation of lattice quantum electrodynamics.
\end{abstract}

\begin{quotation}
	Based on the insight that temporal and spatial instabilities can be 
		treated on equal footing, the `chronotopic' theory of 1990s offers an 
		elegant description of  \spt\ chaos, with predictions of the theory 
		encapsulated in an `entropy density' function%
\rf{Grassberger89,LePoTo96,LePoTo97,LePoTo97a}.
		The first decade of 2000s saw rapid progress in description of 
		transitional turbulence in fluids, in terms of 
%`{\ecs}s', 
		numerically exact {\spt}ly periodic unstable solutions of {\NS} equations.%
	%RESTORE		\rf{W01,WK04,KawKida01,science04,GHCW07,WFSBC15,CPTKGS22}
		This approach to turbulence aims to predict measurable observables from 
		the defining (\NS) equations, without any statistical assumptions. 
		But how are these solutions to be pieced together?%
%RESTORE	\rf{focusPOT}
		As a single such solution requires highly unstable time integration over 
		$10^4-10^6$ discretization ODEs, estimates of `chronotopic' theory 
		observables based on large spacetime volume limits of numerical 
		simulations are out of question. 
\refRef{CL18} develops a new \spt\ theory of turbulence, which
		assigns a new kind of a global weight, an `\Hilldet', to each {\lst},
		with a new enumeration of global solutions as sums over all {\Blatt}s. 
		Our {\dft} replaces numerical extrapolations of finite spacetime 
		simulations of the 1990s chronotopy by our main result, 
		the {exact} \spt\ {\detZeta}, with computable cycle expansion truncations 
		errors, decreasing exponentially with the \spt\ volumes of \lsts\ 
		included in its evaluation.  
\end{quotation}


In 1924 Ryuichiro Kitano\rf{Kitano24} 
% (for details, see \refpage{post:Kitano24})
applied in his
{\em QED Five-Loop on the Lattice},
% Progress of Theoretical and Experimental Physics, Volume 2025, 2025, 013B02, 
%Prog. Theor. Exp. Phys. 2025 013B02 (10 pages) DOI: 10.1093/ptep/ptae194
%\HREF{https://doi.org/10.1093/ptep/ptae194} {DOI}, 
\arXiv{2411.11554},
the lattice QED stochatic quantization method 
proposed in 
\refrefs{KiTaHa21,KitTak23}, 
% \refpage{post:KiTaHa21} and  \refpage{post:KitTak23},
to evaluation of the quenched QED contribution of the {$g-2$} magnetic moment anomaly
to $\alpha^5$th (5-loop) order.
Kitano, Takaura, and Hashimoto's  methodology is an example of
{\em numerical stochastic perturbation theory (NSPT)}.
%, see the notes of \refsect{sect:NSPT}.

Kitano had performed his Monte-Carlo lattice simulations on five lattice 
\HREF{https://en.wikipedia.org/wiki/Hyperrectangle}{`hyperrectangular'} {\pcell}s,
\[
\BravCell{24^3}{48}{0}, 
\BravCell{28^3}{56}{0}, 
\BravCell{32^3}{64}{0}, 
\BravCell{48^3}{96}{0},
\BravCell{64^3}{128}{0}
\,, 
\]
in the notation of \toCLpdf{equation.74}{eq.~(76)},
and obtained values of the 4 and 5-loop contributions 
% listed in table~\refeq{anomalValues}
that - though not competitive with the Feynman diagrams evaluations - were
in the right ball park.


For systems 
characterized by several translational symmetries, one has to take care 
of multiple periodicities; in the language of crystallography,  organize 
	the \po\ sums by corresponding {\em {\Blatt}s}, or, in the language of 
field theory, by the `sum over \emph{geometries}'. In \refsect{s:Blatts} 
we enumerate and construct spacetime geometries
%, or $d=2$ {\Blatt}s  $\LTS{}{}{}$, 
of increasing spacetime periodicities. The classification 
of {\lsts} proceeds in two steps. On the coordinate level, periodicity is 
imposed by the hierarchy of {\Blatt}s of increasing periodicities. 
% dropped reference to (\refsect{s:Blatt}).
On the field-configuration level, the key to the \spt\ \po\ theory is the 
enumeration and  determination of \emph{\primeOrbs}, the basic building 
blocks of \po\ theory.
%RESTORE	 (\refsect{s:orbits}). 

\PC{2026-08-11} 
{Throughout the manuscript: please include with the equations the source literature that you got them from.}


\section{Hypercubic {\Blatt}s}
\label{s:Blatts} 

Whenever we refer here to a `{\Blatt}', we mean one of the infinity of 
`full rank {sub}lattices of the hypercubic lattice'\rf{Cassels59}. 

{\Blatt}
\beq
\lattice = \Big\{\sum_{i=1}^d \cl{i} {\bf a}_i
\;\vert\; \cl{i} \in\mathbb{Z} \Big\}
\ee{BravLatt}
where the $d$ independent integer lattice vectors
\(
\{{\bf a}_1,{\bf a}_2,\cdots,{\bf a}_d\}
\) % ee{integBas}
where the $[d\!\times\!d]$ basis matrix 
\beq
\Bcell{A} = 
\left[
\mathbf{a}_{1}, \mathbf{a}_{2}, \cdots,  \mathbf{a}_{d}
\right]
\ee{dDimBcell}
formed from primitive integer lattice vectors $\{\mathbf{a}_j\}$ 
defines a $d$\dmn\ \emph{\pcell}\rf{AshMer,CoEvBu22}.


A {\pcell} is not unique\rf{Siegel89}: the {\Blatt} 
$\lattice_{\Bcell{A}'}$ defined by basis $\Bcell{A}'$ is the same as the 
{\Blatt} $\lattice_\Bcell{A}$ defined by basis 
\( %beq
\Bcell{A}=\Bcell{A}'\,\Bcell{U} 
\) %ee{UnimdTran}
if the two are related by a $[d\!\times\!d]$ unimodular, volume 
preserving 
% $\Det\Bcell{U}=\pm1$
matrix  $\Bcell{U}\in\SLn{d}{\integers}$ 
transformation\rf{Lang71,Cardy86,ZiLoKl99}.
%see \reffig{f:BravLatt}\,(b). 
Constructing 
\emph{all} {\Blatt}s, however, is straightforward, as each such infinite 
family of equivalent {\pcell}s contains a single, unique \emph{Hermite 
	normal form} {\pcell}, with upper-triangular basis\rf{Cohen93} primitive 
vectors
\beq 
\Bcell{A} = 
\left[\begin{array}{cc} 
	\speriod{} & \tilt{} \cr
	0 & \period{} 
\end{array}\right]
\,,\qquad \vol_\Bcell{A} = \speriod{} \period{}
\,, 
\ee{dHermiteBas} 
where $\speriod{}$, $\period{}$ are the spatial, 
temporal lattice periods, respectively, and $\vol_\Bcell{A}$ is the 
lattice {volume} \refeq{lattVol}. The tilt\rf{OKKH99} 
$0 \leq \tilt{}<\speriod{}$ imposes `relative-periodic shift' 
{\bcs}\rf{ChaosBook}. 

\subsection{Periodic {\fconfs}}
\label{s:periodicCnfs} 

	We shall demand that {\ELe} are the same everywhere in spatially
	translation-invariant directions, for all times. 
	The only way to obey that is by solutions $\Xx_c$  
	being periodic.

A scalar field $\ssp(x)$ over $d$ Euclidean coordinates can be
discretized by
replacing the continuous space by a $d$\dmn\ hypercubic {integer lattice}
$\integers^d$,
%
% dasbuch/videos/Camtasia Studio/diffusion/00diffusionVideos.txt
% Title: Lattice field theory
% LectureW4L2latt.tscproj    1/6        X:22 min
% maybe RESTORE \toVideo{youtube.com/embed/EWLQJ6ZpUWQ}
%
with lattice spacing $a$, and
evaluating the {field} only on the
lattice points\rf{MonMun94,MunWal00}
\beq
\ssp_z
=
\ssp(x)
\,,\qquad \qquad x = az= \mbox{lattice point}
\,,\quad z \in \integers^d
\,.
\ee{LattField}
% see \reffig{FieldConfig}.
A {\em field configuration} (here in one {\spt} dimension)
\beq
\Xx =
\cdots {\ssp}_{-3} {\ssp}_{-2}\,{\ssp}_{-1}\,
{\ssp}_0\,
{\ssp}_{1} {\ssp}_{2} {\ssp}_{3} {\ssp}_{4}  \cdots
\,,
\ee{stateSp}
takes any set of values $\ssp_{z}\in\reals$ in system's $\infty$\dmn\
\emph{\statesp}.
A {periodic} {field configuration} $\Xx$ satisfies
\beq
\Xx({x} + {R}) = \Xx({x})
\ee{1DprimePO}
for any discrete translation
$R$ in the the {\Blatt} \eqref{BravLatt}.

We say that a {lattice {\fconf}} is $\lattice_\Bcell{A}$-\emph{periodic} if 
\beq 
%\Xx({x} + {R}) = \Xx({x})
\ssp_{z+\shift} = \ssp_z 
\ee{periodLConf}  % was {1DprimePO}
for any discrete translation 
$\shift=n_{1}\mathbf{a}_1+n_{2}\mathbf{a}_2+\cdots+n_{d}\mathbf{a}_d$ in 
the {\Blatt} \eqref{BravLatt}
where the $[d\!\times\!d]$ basis matrix 
\(
\Bcell{A} = 
\left[
\mathbf{a}_{1}, \mathbf{a}_{2}, \cdots,  \mathbf{a}_{d}
\right]
\) 
formed from primitive integer lattice vectors $\{\mathbf{a}_j\}$ 
defines a $d$\dmn\ \emph{\pcell}\rf{AshMer,CoEvBu22}.
If the lattice 
spacing \refeq{LattField} is set to 1, the lattice \emph{volume}, or the 
volume of the {\pcell} 
\beq
\vol_\Bcell{A} = |\Det\Bcell{A}| 
\ee{lattVol} 
equals the number of lattice sites $z\in \Bcell{A}$ within the {\pcell}.
%,  see \reffig{f:BravLatt}. 

\subsection{Integer {\Blatt}s}
\label{s:Blatt1d} 


\subsection{Square lattice {\Blatt}s}
\label{s:BlattSquare} 

Consider a $[2\!\times\!2]$ integer basis matrix \refeq{BravLatt} \beq 
\Bcell{A} = \left[ \mathbf{a}_{1}, \mathbf{a}_{2} \right] 
= \left[
\begin{array}{cc}
	{a}_{1,1} & {a}_{2,1}\cr
	{a}_{1,2} & {a}_{2,2}
\end{array}
\right]
\,,\quad
\mathbf{a}_{j} = \left[ \begin{array}{c} {a}_{j,1} \cr {a}_{j,2} 
\end{array} \right]
\,, 
\ee{2dPrimCell} 
formed from a pair of two\dmn\ integer lattice 
primitive vectors $\mathbf{a}_{1}$, $\mathbf{a}_{2}$. 
% as in \reffig{f:BravLatt}.
A two\dmn\ {\em {\Blatt}}, 
% \reffig{f:BravLatt}, 
\beq 
\lattice_\Bcell{A} = 
\left\{\Bcell{A} \mathbf{n}  \,|\,\mathbf{n} \in \mathbb{Z}^2 \right\} 
%\,,
\ee{2DBravaisLattice} 
%of lattice points
generated by all discrete translations 
\( 
% \shift =
\Bcell{A}  \mathbf{n} 
% = m_{1}\mathbf{a}_{1}+m_{2}\mathbf{a}_{2}
%\,.
\) 
is a sublattice of the integer lattice $\integers^2$. 

As in a discretized field theory the fields are defined only on the 
hypercubic integer lattice, not on a continuum, we define the {\em	\pcell} 
%\refeq{torusStatesp} 
as the set of lattice sites within the 
parallelepiped \refeq{2dPrimCell}.
% illustrated by \reffig{f:BravLatt}. 
The  tips of primitive vectors and parallelepiped's outer boundaries belong, 
by translation, to the neighboring tiles; this yields the correct lattice 
{volume} \refeq{lattVol}, the number of lattice sites $\vol_\Bcell{A}$ 
within the {\pcell} $\Bcell{A}$. 

% formed by the primitive vectors $\Bcell{A}$,

Constructing  square lattice 
\emph{all} {\Blatt}s is straightforward, as each  infinite 
family of equivalent {\pcell}s contains a single, unique \emph{Hermite 
	normal form} {\pcell}, with upper-triangular basis\rf{Cohen93} primitive 
vectors $\mathbf{a}_1=(\speriod{}, 0)$, 
$\mathbf{a}_2=(\tilt{},\period{})$, 
\beq 
\Bcell{A} = 
\left[\begin{array}{cc} 
	\speriod{} & \tilt{} \cr
	0 & \period{} 
\end{array}\right]
\,,\qquad \vol_\Bcell{A} = \speriod{} \period{}
\,, 
\ee{2DHermiteBas} 
where $\speriod{}$, $\period{}$ are the spatial, 
temporal lattice periods, respectively, and $\vol_\Bcell{A}$ is the 
lattice {volume} \refeq{lattVol}. The tilt\rf{OKKH99} 
$0 \leq \tilt{}<\speriod{}$ imposes `relative-periodic shift' 
{\bcs}\rf{ChaosBook}. 

This equivalence underlies many of the properties of elliptic functions 
and modular forms\rf{SteSha03} (see \refeq{2tauBas}). 
	In the theory of elliptic functions\rf{SteSha03} the {\pcell} is 
	represented by a complex modular parameter $\tau$, with spatial period 
	$\speriod{}$ taken as the lattice spacing constant \refeq{LattField}, 
	primitive vectors $\mathbf{a}_1=(1, 0)$, $\mathbf{a}_2=(\tau_1,\tau_2)$, 
	so $\period{}\to\tau_2=\period{}/\speriod{}$, 
	$\tilt{}\to\tau_1=\tilt{}/\speriod{}$, and 
	\beq 
	\Bcell{A} = 
	\left[\begin{array}{cc} 
		1 & \tau_1 \cr
		0 & \tau_2 
	\end{array}\right]
	\,,\qquad 
	%\vol_\Bcell{A}  = 
	|\Det\Bcell{A}|= \tau_2 
	\,, 
	\ee{2tauBas} 
	`Hermite normal form' 
	corresponds to  the modular parameter $\tau$ values in the {fundamental 
		domain}. If the corresponding torus is visualized as glueing of a unit 
	square into a tube, $\tau_2$ parameterizes how the tube is stretched, and 
	$\tau_1$ parameterizes how it is twisted before its edges are stitched 
	together. 
	
	Here we refer to a particular {\Blatt} by its Hermite normal form 
	\refeq{2DHermiteBas}, as 
	\beq 
	\lattice_\Bcell{A}=\LTS{}{}{} 
	\,, 
	\ee{LxT_S} 
	and to the set of lattice sites within the primitive 
	parallelogram \Bcell{A} as its {\pcell}. Notation $\LTS{}{}{}$ refers to 
	{\pcell} being a rectangle of spatial width $\speriod{}$, temporal height 
	$\period{}$, with the {\pcell} above it shifted by $S$, see, for example 
	the $\BravCell{3}{2}{1}$ {\pcell} in
	%shown in \reffig{f:2x1rpo}\,(b). In 
	terms of lattice site fields, a {\fconf} $\ssp_{z_1 z_2}$ 
	\refeq{periodLConf}, $z_1z_2\in\integers^2$, satisfies the 
	$\tilt{}$-corkscrew boundary condition\rf{ChMaVa96}, 
	\bea 
	\mbox{horizontally: } \quad && \ssp_{z_1 z_2} \,=\, \ssp_{z_1 + 
		\speriod{}, z_2} 
	\continue
	\mbox{vertically: }   \quad && \ssp_{z_1 z_2} \,=\, \ssp_{z_1 + \tilt{}, 
		z_2 + \period{}} 
	\,.
	\label{screw} 
	\eea 
	%see \reffig{f:2x1rpo}. 
	
\subsection{Cubic lattice {\Blatt}s}
\label{s:Blatt3d2} 

\subsection{Four\dmn\ hypercubic {\Blatt}s}
\label{s:Blatt4d} 


\section{The  U(1) lattice action and \ELe}

We consider  $U(1)$ lattice electrodynamics formulated on a four-dimensional Euclidean Bravais lattice $\lattice_{\Bcell{A}} = [L_1 \times L_2 \times L_3 \times T]_S$ subject to periodic boundary conditions along all coordinate directions. The lattice sites are designated by integer coordinates $n = (n_1, n_2, n_3, n_4)$. The gauge fields are parameterized by directed link variables
\begin{equation}
U_\mu(n) = \exp\left(i e A_\mu(n)\right),
\end{equation}
where $e$ denotes the bare electromagnetic coupling constant and $A_\mu(n) \in \mathbb{R}$ represents the angular gauge potential associated with the link emanating from site $n$ in the positive unit direction $\hat{\mu}$ for $\mu \in \{1,2,3,4\}$.

\PC{2026-08-11}
{Why are you using "angular gauge potential"?}

The Euclidean dynamics of the pure gauge sector is governed by the Wilson plaquette action augmented by an infrared-regulating photon mass term:
\begin{equation}
S_{\Bcell{A}}[A] = \beta \sum_{n \in \Bcell{A}} \sum_{1 \le \mu < \nu \le 4} \left[ 1 - \cos\left(\theta_{\mu\nu}(n)\right) \right] + \frac{1}{2} m_\gamma^2 \sum_{n \in \Bcell{A}} \sum_{\mu=1}^4 A_\mu(n)^2.
\end{equation}
Here, $\beta = 1/e^2$ is the dimensionless lattice gauge coupling, $m_\gamma$ denotes the infrared photon mass regulator, and $\theta_{\mu\nu}(n)$ is the gauge-invariant, oriented field-strength angle around an elementary lattice plaquette in the $(\mu,\nu)$ plane:
\begin{equation}
\theta_{\mu\nu}(n) = e \left( A_\mu(n) + A_\nu(n+\hat{\mu}) - A_\mu(n+\hat{\nu}) - A_\nu(n) \right).
\end{equation}

The deterministic classical configurations of the lattice theory correspond to the stationary points of the Euclidean functional. Requiring the first variation of the action with respect to every link potential to vanish identically, $\delta S_{\Bcell{A}} / \delta A_\mu(n) = 0$, yields a coupled system of nonlinear difference equations:
\begin{equation}
\sum_{\nu \neq \mu} \left[ \sin\left(\theta_{\mu\nu}(n)\right) - \sin\left(\theta_{\mu\nu}(n-\hat{\nu})\right) \right] + m_\gamma^2 e^{-2} A_\mu(n) = 0.
\end{equation}
This system represents the exact discrete Euler-Lagrange equations of motion. Because of the periodic dependence on the field-strength tensor, the configuration space supports multiple distinct stationary solutions.

\section{Algebraic reduction via maximal spanning tree gauge}

When $m_\gamma \to 0$, the action $S_{\Bcell{A}}$ remains invariant under local site-dependent gauge transformations $A_\mu(n) \to A_\mu(n) + e^{-1}[\alpha(n+\hat{\mu}) - \alpha(n)]$. In numerical root-finding protocols, continuous gauge degeneracies produce zero-eigenvalue directions in the {\jacobianOrb}, rendering the Jacobian singular. Differential gauge-fixing conditions enforce constraints through Lagrange multipliers, which introduce saddle-point indefiniteness and Gribov copies into the stationary spectrum.

To ensure a positive-definite and non-singular Jacobian around classical minima, we eliminate local gauge redundancy algebraically by imposing the Maximal Spanning Tree gauge. On a finite lattice containing $V$ vertices, a maximal spanning tree $T_{\text{max}}$ is a connected subgraph containing all $V$ vertices and exactly $V-1$ links without closed cycles. By transforming the link variables along $T_{\text{max}}$ to the identity, we set exactly $V-1$ gauge potentials to zero:
\begin{equation}
A_\mu(n) = 0 \quad \text{for all } (n, \hat{\mu}) \in T_{\text{max}}.
\end{equation}

We restrict our computational analysis to the minimal Euclidean primitive cell $[2^3 \times 2]_0$, where each coordinate $n_\mu \in \{0, 1\}$ is defined modulo 2. This toroidal lattice contains $V = 16$ vertices and a total of $4V = 64$ directed link variables. We designate the origin $(0,0,0,0)$ as the root vertex and construct an explicit breadth-first maximal spanning tree:
\begin{equation}
\begin{split}
T_{\text{max}} = {}& \left\{ (n, \hat{1}) \mid n_1 = 0 \right\} \cup \left\{ (n, \hat{2}) \mid n_1 \in \{0,1\}, n_2 = 0 \right\} \\
& \cup \left\{ (n, \hat{3}) \mid n_1,n_2 \in \{0,1\}, n_3 = 0 \right\} \cup \left\{ (n, \hat{4}) \mid n_1,n_2,n_3 \in \{0,1\}, n_4 = 0 \right\}.
\end{split}
\end{equation}

Imposing $A_\mu(n) = 0$ along these $V-1 = 15$ tree links removes all local gauge freedom. The physical configuration space is thereby projected onto the remaining $64 - 15 = 49$ active dynamical link variables, which include the homologically non-trivial Polyakov loops wrapping around the four toroidal boundaries.

\section{Construction of {\primeOrb}s}

Within the pure gauge limit ($m_\gamma = 0$), any field configuration whose oriented plaquette angles $\theta_{\mu\nu}(n)$ are strictly uniform across adjacent lattice sites satisfies the Euler-Lagrange equations identically, as the spatial difference of sines vanishes at every vertex. We derive the explicit classical action $S_{\Bcell{A}}$ for the four fundamental classes of stationary prime orbits supported on the $[2^3 \times 2]_0$ primitive cell.

\subsection{The trivial vacuum orbit}
\PC{2026-08-16}{"trivial vacuum orbit" you mean zero field? Or \steady?}
The absolute minimum of the Euclidean action is given by the homogeneous zero-field configuration where all 49 active link variables vanish:
\begin{equation}
A_\mu(n) = 0 \quad \text{for all active links.}
\end{equation}
Every elementary plaquette angle equals zero, $\theta_{\mu\nu}(n) = 0$, which yields $\sin(0) = 0$ identically across the lattice. The classical action of this ground-state reference orbit is
\begin{equation}
S^{(0)} = 0.
\end{equation}

\subsection{Constant-flux toron orbits}
The toroidal geometry of the primitive cell supports topologically non-trivial winding modes carrying quantized homogeneous flux through any of the six coordinate planes $(\mu,\nu)$. Let $k_{\mu\nu} \in \mathbb{Z}$ represent the integer flux quantum threading the $(\mu,\nu)$ plane. Because every coordinate plane on the $[2^3 \times 2]_0$ lattice has a two-dimensional area of $L_\mu L_\nu = 4$, a single flux quantum ($k_{\mu\nu} = 1$) requires a total integrated phase of $2\pi$ distributed uniformly over four elementary plaquettes:
\begin{equation}
\theta_{\mu\nu}(n) = \frac{2\pi k_{\mu\nu}}{L_\mu L_\nu} = \frac{\pi}{2}.
\end{equation}
For the orientation plane $(1,2)$, this stationary state is realized by assigning boundary shear potentials to the active links wrapping the periodic borders:
\begin{equation}
A_2(1, n_2, n_3, n_4) = \frac{\pi}{2 e}, \quad A_1(n_1, 1, n_3, n_4) = -\frac{\pi}{2 e},
\end{equation}
with all remaining active link variables set to zero. Because $\sin(\pi/2) = 1$ is invariant across all sixteen vertices, the discrete divergence of the field-strength tensor vanishes identically. Summing the Wilson plaquette energy over the sixteen lattice sites for a single excited coordinate plane yields the classical action:
\begin{equation}
S_{\text{toron}} = \beta \sum_{n \in \Bcell{A}} \left[ 1 - \cos\left(\frac{\pi}{2}\right) \right] = 16\beta.
\end{equation}
There exist six independent fundamental toron orbits corresponding to the six orientation planes of the four-dimensional torus.

\subsection{Saddle-point sphaleron orbits}
Sphalerons are anharmonic, unstable stationary states where the field-strength angles of an entire coordinate plane are excited to the apex of the periodic cosine potential barrier:
\begin{equation}
\theta_{\mu\nu}(n) = \pi.
\end{equation}
Because $\sin(\pi) = 0$, the discrete Euler-Lagrange equations are satisfied at every vertex. For the $(1,2)$ coordinate plane, this saddle-point configuration is constructed via half-turn boundary windings:
\begin{equation}
A_2(1, n_2, n_3, n_4) = \frac{\pi}{e}, \quad A_1(n_1, 1, n_3, n_4) = 0.
\end{equation}
Evaluating the Wilson action across the sixteen sites gives:
\begin{equation}
S_{\text{sphaleron}} = \beta \sum_{n \in \Bcell{A}} \left[ 1 - \cos(\pi) \right] = 32\beta.
\end{equation}
These six sphaleron configurations act as discrete separating boundaries between distinct topological vacuum sectors.

\subsection{Monopole-antimonopole dipole orbits}
Compact $U(1)$ lattice gauge theory supports localized topological excitations in the form of DeGrand-Toussaint magnetic monopoles and Dirac strings. A magnetic monopole is identified within a three-dimensional elementary cube when the oriented sum of the six surrounding plaquette angles satisfies
\begin{equation}
\frac{1}{2\pi} \sum_{\text{cube}} \bar{\theta}_{\mu\nu} = m \in \mathbb{Z} \setminus \{0\}.
\end{equation}
Because the total magnetic charge on a closed toroidal manifold must vanish, localized topological defects occur exclusively as neutral monopole-antimonopole dipole pairs. On the $[2^3 \times 2]_0$ primitive cell, the minimal-energy configuration places a monopole at the origin $(0,0,0,0)$ and an antimonopole at the maximally separated antipodal site $(1,1,1,0)$. Numerical relaxation of the 49 active link variables converges to a localized standing wave where the core plaquettes take stationary anharmonic values concentrated near $\theta_{\mu\nu} \approx \pm \pi/3$ and $\pm 2\pi/3$. The numerical evaluation of the Wilson action for this antipodal defect geometry yields:
\begin{equation}
S_{\text{dipole}} \approx 17.83\beta.
\end{equation}
There exist four distinct antipodal dipole orientations corresponding to alignment along the four coordinate axes.

\section{Linear stability analysis and {Morse} indices}

To determine the contribution of each prime orbit to the deterministic spatiotemporal zeta function, we analyze the local linear stability spectrum around each stationary solution. We define the $49 \times 49$ reduced Hessian stability matrix restricted to the active tree-gauge link variables:
\begin{equation}
H_{ij} = \frac{\partial^2 S_{\Bcell{A}}}{\partial A_i \partial A_j}.
\end{equation}
The off-diagonal and diagonal entries of $H_{ij}$ are weighted by local curvature factors of $\beta \cos(\theta_{\mu\nu}(n))$. The Morse index $M$ of a stationary orbit is defined as the number of strictly negative eigenvalues of $H_{ij}$, which corresponds to the number of unstable hyperbolic directions in configuration space.

For the trivial vacuum orbit, $\cos(0) = 1$ across all 96 elementary plaquettes of the lattice. The matrix $H_{ij}$ coincides with the positive-definite discrete Laplace-Beltrami operator on the 49-dimensional active subspace. All 49 eigenvalues are strictly positive ($\lambda_k > 0$), establishing a Morse index of $M = 0$.

For the constant-flux toron orbit, $\cos(\pi/2) = 0$ along the 16 plaquettes of the excited coordinate plane, while $\cos(0) = 1$ on the remaining 80 unexcited plaquettes. This vanishing curvature along the flux direction produces zero-eigenvalues that correspond to marginal toroidal flux translations, while all remaining transverse eigenvalues remain strictly positive. Consequently, the toron possesses no negative eigenvalues, yielding a Morse index of $M = 0$.

For the sphaleron orbit, $\cos(\pi) = -1$ along all 16 plaquettes of the excited plane. This negative curvature injects negative eigenvalues into the reduced {\jacobianOrb}. Diagonalizing $H_{ij}$ reveals exactly 16 negative eigenvalues corresponding to the independent directions along which the plaquette angles can roll off the potential barrier, establishing a Morse index of $M = 16$.

For the antipodal monopole-antimonopole dipole orbit, regions of positive curvature ($\cos \theta > 0$) far from the core surround regions of negative curvature ($\cos \theta < 0$) localized along the defect core and the connecting Dirac string. Spectral factorization of the $49 \times 49$ {\jacobianOrb} evaluated at the stationary dipole configuration reveals exactly 4 negative eigenvalues, yielding a Morse index of $M = 4$.

\section{Synthesis of the deterministic zeta function weight catalog}

In the spatiotemporal dynamical systems formalism, the contribution of a classical prime orbit $p$ to the deterministic zeta function is governed by its statistical weight factor:
\begin{equation}
w_p = (-1)^{M_p} \exp\left(-S_{\Bcell{A}}^{(p)}\right).
\end{equation}
The phase factor $(-1)^{M_p}$ accounts for the Morse index of the hyperbolic trajectory, ensuring the correct sign parity in the periodic-state trace formula.

We assemble the complete operational catalog of fundamental prime orbits derived on the minimal primitive cell $[2^3 \times 2]_0$ in Table \ref{tab:zeta_catalog}.

\begin{table}[h!]
\centering
\caption{Catalog of fundamental prime orbits on the $\Bcell{A}=[2^3 \times 2]_0$ {\pcell} $U(1)$ lattice, detailing topological multiplicity,  Euclidean action, Morse index, and deterministic zeta function weight.}
\label{tab:zeta_catalog}
\begin{tabular}{@{}lcccc@{}}
\toprule
\shortstack[l]{Orbit\\Classification} & \shortstack[c]{Topological\\Multiplicity} & \shortstack[c]{Classical Action\\$S_{\text{compact}}$} & \shortstack[c]{Morse\\Index $M$} & \shortstack[c]{Zeta Function\\Weight $w_p$} \\
\midrule
Trivial Vacuum & 1 & $0$ & 0 & $+1$ \\
Constant-Flux Torons & 6 & $16\beta$ & 0 & $+6 \exp(-16\beta)$ \\
Antipodal Dipole Pairs & 4 & $17.83\beta$ & 4 & $+4 \exp(-17.83\beta)$ \\
Saddle-Point Sphalerons & 6 & $32\beta$ & 16 & $+6 \exp(-32\beta)$ \\
\bottomrule
\end{tabular}
\end{table}

This closed catalog provides the exact non-perturbative weights required for the initial phase of the deterministic evaluation of lattice quantum electrodynamics. By substituting these classical action values and stability spectra into the periodic-state averaging formula, one establishes a convergent mathematical foundation for evaluating lattice observables without relying on stochastic Monte-Carlo sampling.

%  Too busy for Predrag :)
%\begin{figure*}
%    \centering
%    \includegraphics[width=0.8\linewidth]{illustration.png}
%    \captionsetup{labelformat=empty}
%    \caption{Workflow}
%    \label{fig:enter-label}
%\end{figure*}

%%%%%%%%%%%%%%%%%%%%%%%%%%%%%%%%%%%%%%%%%%%%%%%%
\printbibliography[heading=subbibintoc,title={References}]

\end{document}